\documentclass[10pt, aspectratio=169]{beamer}
\input{../../_en_preambulo_beamer}
% Comandos y macros estadísticas compartidas para las presentaciones Beamer en inglés (Metropolis)

% Conjuntos numéricos básicos
\newcommand{\R}{\mathbb{R}}
\newcommand{\N}{\mathbb{N}}
\newcommand{\Q}{\mathbb{Q}}
\newcommand{\Z}{\mathbb{Z}}
\newcommand{\set}[1]{\left\{ #1 \right\}}
\newcommand{\abs}[1]{\left|#1\right|}
\newcommand{\norm}[1]{\left\| #1 \right\|}

% Comandos estadísticos y probabilísticos del curso
\newcommand{\Var}{\operatorname{Var}}
\newcommand{\cov}{\operatorname{Cov}}
\newcommand{\corr}{\rho}
\newcommand{\comb}[2]{\begin{pmatrix} #1 \\ #2 \end{pmatrix}}
\newcommand{\s}{\sigma}
\newcommand{\particion}{\mathfrak{P}}
\newcommand{\card}[1]{\#\left( #1 \right)}
\newcommand{\seq}[3]{\set{#1}_{#2}^{#3}}
\newcommand{\E}{\mathbb{E}}
\newcommand{\Prob}{\mathbb{P}}

% Entornos pedagógicos y cajas en inglés académico natural (soportando tanto nombres en inglés como en español)
\newenvironment{discussion}{%
  \begin{alertblock}{Discussion Question}%
}{%
  \end{alertblock}%
}
\newenvironment{discusion}{%
  \begin{alertblock}{Discussion Question}%
}{%
  \end{alertblock}%
}

\newenvironment{reflection}{%
  \begin{alertblock}{Classroom Reflection}%
}{%
  \end{alertblock}%
}
\newenvironment{reflexion}{%
  \begin{alertblock}{Classroom Reflection}%
}{%
  \end{alertblock}%
}

\newenvironment{paradox}{%
  \begin{alertblock}{Exploratory Paradox}%
}{%
  \end{alertblock}%
}
\newenvironment{paradoja}{%
  \begin{alertblock}{Exploratory Paradox}%
}{%
  \end{alertblock}%
}

\newenvironment{motivation}{%
  \begin{exampleblock}{Motivation: Data Science in Practice}%
}{%
  \end{exampleblock}%
}
\newenvironment{motivacion}{%
  \begin{exampleblock}{Motivation: Data Science in Practice}%
}{%
  \end{exampleblock}%
}

\newenvironment{quickchallenge}{%
  \begin{block}{Quick Team Challenge}%
}{%
  \end{block}%
}
\newenvironment{retoclase}{%
  \begin{block}{Quick Team Challenge}%
}{%
  \end{block}%
}


\title[Nonlinear Transformations]{Section 09.12: Nonlinear Transformations and Polynomial Regression}
\subtitle{Closing the Chapter: Curvature, Polynomial Multicollinearity, and Synthesis}
\author[J. Castillo Colmenares]{Juliho Castillo Colmenares}
\institute{Tecnológico de Monterrey}
\date{\vspace{-1.2cm}}

\begin{document}

% Slide 1: Title
\begin{frame}[plain]
  \titlepage
\end{frame}

% Slide 2: Roadmap
\begin{frame}{Roadmap --- Chapter 09: Linear and Multiple Regressions}
  \scriptsize
  Where are we in the chapter's modular development?
  \begin{itemize}\itemsep=0.02em
    \item \textbf{09.01} Correlation as a Premise for Regression
    \item \textbf{09.02} Introduction to Linear Regression
    \item \textbf{09.03} Mathematics of Regression: OLS Derivation and $R^2$
    \item \textbf{09.04} Linear Regression on Simulated Data
    \item \textbf{09.05} Optimal Coefficients, $t$/$F$ Tests, and RSE
    \item \textbf{09.06} Implementation in Python with \texttt{statsmodels}
    \item \textbf{09.07} Multiple Regression, the Normal Equation, and Ridge/Lasso
    \item \textbf{09.08} Model Validation and $k$-fold Cross-Validation
    \item \textbf{09.09} Residual Diagnostics and Multicollinearity (VIF)
    \item \textbf{09.10} Regression with \texttt{scikit-learn}
    \item \textbf{09.11} Categorical and Dummy Variables
    \item \textbf{09.12} Nonlinear Transformations and Polynomial Regression \emph{(Today --- Chapter Closing)}
  \end{itemize}
  \pause
  \vspace{-0.05cm}
  \begin{block}{Session Objective}
    Model curved relationships (nonlinear in the original variable) by reformulating them as multiple linear regression on expanded variables; diagnose when a transformation is needed; and close Chapter 09 by synthesizing its full arc, from correlation to polynomial regression.
  \end{block}
\end{frame}

% Slide 3: Motivation
\begin{frame}{When the Relationship Is Not a Straight Line}
  \footnotesize
  \begin{motivacion}
    Engine horsepower (\texttt{hp}) and fuel efficiency (\texttt{mpg}) do not always follow a simple monotonic trend: they can trace out a curve. A strictly linear model $\texttt{mpg}=c_0+c_1\,\texttt{hp}$ cannot represent a change of direction in the curve.
  \end{motivacion}

  \vspace{0.15cm}
  \pause
  \begin{itemize}\itemsep=0.1cm
    \item \textbf{Diagnosis:} plotting the response against each predictor reveals the true shape of the relationship. \pause
    \item \textbf{Solution:} add higher-order terms ($X^2$, $X^3$, \ldots) or apply transformations (logarithm, square root).
  \end{itemize}
\end{frame}

% Slide 4: Polynomial regression as a case of multiple regression
\begin{frame}{Polynomial Regression Is Multiple Regression}
  \footnotesize
  \begin{block}{The Auxiliary-Variable Trick}
    By defining $Z=X^2$ as a new, explicitly computed column, the model $Y=c_0+c_1X+c_2X^2+\varepsilon$ is rewritten exactly as $Y=c_0+c_1X+c_2Z+\varepsilon$.
  \end{block}
  \pause

  \vspace{0.1cm}
  \begin{alertblock}{The Normal Equation Does Not Change}
    Least squares only requires linearity \textbf{in the parameters} $c_0,c_1,c_2$, not in $X$. The Normal Equation from Section 09.07, $\hat{\mathbf{c}}=(\mathbf{X}^T\mathbf{X})^{-1}\mathbf{X}^T\mathbf{Y}$, applies without any modification.
  \end{alertblock}
\end{frame}

% Slide 5: Polynomial multicollinearity
\begin{frame}{The Inherent Multicollinearity of Polynomials}
  \footnotesize
  \begin{block}{The Problem}
    When $X$ is confined to a narrow range, $X^2\approx X_0^2+2X_0(X-X_0)$ is nearly a linear function of $X$: $X$ and $X^2$ become highly correlated, inflating their VIF (Section 09.09).
  \end{block}
  \pause

  \vspace{0.1cm}
  \begin{alertblock}{It Does Not Invalidate Predictions}
    Predictions $\hat Y$ (which depend on the joint combination of $X$ and $X^2$) remain accurate; what becomes difficult is interpreting $c_1$ and $c_2$ in isolation.
  \end{alertblock}
\end{frame}

% Slide 6: Python Lab (Block 1)
\begin{frame}[fragile]{Python Lab: Fitting the Quadratic Model}
  \scriptsize
  Expanding the design matrix and solving via the Normal Equation (\texttt{09.12\_nonlinear\_polynomial\_regression.py}):
  {\lstset{basicstyle=\fontsize{3.2pt}{3.9pt}\selectfont\ttfamily,aboveskip=0.05em,belowskip=0.05em}
  \lstinputlisting[firstline=17, lastline=33]{../../code/09_regresiones/09.12_nonlinear_polynomial_regression.py}}
\end{frame}

% Slide 7: Python Lab (Block 2)
\begin{frame}[fragile]{Python Lab: Linear vs. Quadratic}
  \scriptsize
  Comparing $R^2$ and the Partial F-Test on the quadratic term:
  {\lstset{basicstyle=\fontsize{3.0pt}{3.7pt}\selectfont\ttfamily,aboveskip=0.05em,belowskip=0.05em}
  \lstinputlisting[firstline=38, lastline=59]{../../code/09_regresiones/09.12_nonlinear_polynomial_regression.py}}
\end{frame}

% Slide 8: Python Lab (Block 3)
\begin{frame}[fragile]{Python Lab: Polynomial Multicollinearity}
  \scriptsize
  Correlation between $X$ and $X^2$ depending on the range of $X$:
  {\lstset{basicstyle=\fontsize{3.2pt}{3.9pt}\selectfont\ttfamily,aboveskip=0.05em,belowskip=0.05em}
  \lstinputlisting[firstline=62, lastline=75]{../../code/09_regresiones/09.12_nonlinear_polynomial_regression.py}}
\end{frame}

% Slide 9: Terminal Output
\begin{frame}[fragile]{Python Lab: Terminal Output}
  \scriptsize
  Standard output generated by running the lab script:
  \vspace{0.02cm}
  \begin{block}{Standard Console Output}
    \fontsize{3.2pt}{3.9pt}\selectfont
    \begin{verbatim}
=== Block 1: Quadratic Fit ===
Coefficients (c0,c1,c2): [64.2, -0.7634, 0.003657]; SSE=0.4571

=== Block 2: Linear vs. Quadratic ===
R^2 linear=0.0800, R^2 quadratic=0.9943, Partial F=320.0000

=== Block 3: Polynomial Multicollinearity ===
Narrow X in [95,105]: corr(X,X^2)=0.9999, VIF=5250.82
Wide X in [1,200]:    corr(X,X^2)=0.9673, VIF=15.56
    \end{verbatim}
  \end{block}
\end{frame}

% Slide 10A: Exercise Level 1 (Statement)
\begin{frame}{In-Class Exercise --- Fundamental Level (1/2)}
  \footnotesize
  \begin{block}{Problem (Statement, Problem 9.20.1)}
    A scatter plot of \texttt{hp} vs. \texttt{mpg} shows a descending curve that flattens out and slightly rebounds (an asymmetric inverted-U shape). Identify the appropriate transformation/specification and explain why a linear model would be inadequate.
  \end{block}

  \vspace{0.15cm}
  \pause
  \begin{alertblock}{Strategy}
    A linear model can only capture monotonic trends; consider adding an $X^2$ term.
  \end{alertblock}
\end{frame}

% Slide 10B: Exercise Level 1 (Solution)
\begin{frame}{In-Class Exercise --- Fundamental Level (2/2)}
  \scriptsize
  We analyze the described geometric shape: \pause

  \vspace{0.05cm}
  \begin{block}{Diagnosis}
    The non-monotonic pattern (a change of direction) is characteristic of a \textbf{quadratic} relationship. The appropriate specification is $\texttt{mpg}=c_0+c_1\,\texttt{hp}+c_2\,\texttt{hp}^2$.
  \end{block}

  \vspace{0.05cm}
  \pause
  \begin{alertblock}{Interpretation}
    A linear model forces a monotonic trend onto data that changes direction, producing systematically biased residuals and drastically underestimating the fit quality actually achievable.
  \end{alertblock}
\end{frame}

% Slide 11A: Exercise Level 2 (Statement)
\begin{frame}{In-Class Exercise --- Operational Level (1/2)}
  \footnotesize
  \begin{block}{Problem --- Manual Quadratic Fit (Statement, Problem 9.20.4)}
    Given $\texttt{hp}=\{50,75,100,125,150\}$, $\texttt{mpg}=\{35,28,24,26,32\}$, fit $\texttt{mpg}=c_0+c_1\,\texttt{hp}+c_2\,\texttt{hp}^2$ by building $\mathbf{X}=[\mathbf{1},\texttt{hp},\texttt{hp}^2]$ and solving the Normal Equation.
  \end{block}

  \vspace{0.1cm}
  \pause
  \begin{alertblock}{Strategy}
    $\hat{\mathbf{c}}=(\mathbf{X}^T\mathbf{X})^{-1}\mathbf{X}^T\mathbf{Y}$, exactly as in Section 09.07.
  \end{alertblock}
\end{frame}

% Slide 11B: Exercise Level 2 (Solution)
\begin{frame}{In-Class Exercise --- Operational Level (2/2)}
  \scriptsize
  We solve the $3\times3$ Normal Equation system: \pause

  \vspace{0.05cm}
  \begin{block}{Calculation}
    $$\hat c_0\approx64.20, \qquad \hat c_1\approx-0.7634, \qquad \hat c_2\approx0.003657$$
  \end{block}

  \vspace{0.05cm}
  \pause
  \begin{alertblock}{Interpretation}
    The vertex of the parabola ($\texttt{hp}^*=-c_1/(2c_2)\approx104.4$) marks the point of minimum predicted \texttt{mpg} according to the model.
  \end{alertblock}
\end{frame}

% Slide 12A: Exercise Level 3 (Statement)
\begin{frame}{In-Class Exercise --- Analytical Level (1/2)}
  \footnotesize
  \begin{block}{Problem --- Polynomial as Linear Regression (Statement, Problem 9.20.7)}
    Show that fitting $Y=c_0+c_1X+c_2X^2+\varepsilon$ by least squares is formally identical to the multiple regression of Section 09.07 by defining $Z=X^2$, and explain why the Normal Equation applies without modification.
  \end{block}

  \vspace{0.1cm}
  \pause
  \begin{alertblock}{Strategy}
    Distinguish between linearity in the parameters ($c_0,c_1,c_2$) and linearity in the original predictor ($X$).
  \end{alertblock}
\end{frame}

% Slide 12B: Exercise Level 3 (Solution)
\begin{frame}{In-Class Exercise --- Analytical Level (2/2)}
  \scriptsize
  We rewrite the model with the auxiliary variable: \pause

  \vspace{0.05cm}
  \begin{block}{Calculation}
    \scriptsize
    $$Y=c_0+c_1X+c_2X^2+\varepsilon \quad \xrightarrow{Z=X^2} \quad Y=c_0+c_1X+c_2Z+\varepsilon \quad \text{(linear in } c_0,c_1,c_2\text{)}$$
  \end{block}

  \vspace{0.05cm}
  \pause
  \begin{alertblock}{Interpretation}
    Least squares only requires linearity in the \textbf{parameters}, never in the original predictor variables. That is why the Normal Equation $(\mathbf{X}^T\mathbf{X})^{-1}\mathbf{X}^T\mathbf{Y}$ solves the problem without any further generalization.
  \end{alertblock}
\end{frame}

% Slide 13A: Exercise Level 4 (Statement)
\begin{frame}{In-Class Exercise --- Challenging Level (1/2)}
  \footnotesize
  \begin{block}{Problem --- Polynomial Multicollinearity (Statement, Problem 9.20.9)}
    For $X\in[95,105]$ (a narrow range), show that $X$ and $X^2$ are highly correlated, and explain --- connecting to Section 09.09 --- why this inflates the VIF without invalidating predictions.
  \end{block}

  \vspace{0.1cm}
  \pause
  \begin{alertblock}{Strategy}
    Write $X=X_0+\delta$ with $|\delta|$ small and expand $X^2$ in a Taylor series.
  \end{alertblock}
\end{frame}

% Slide 13B: Exercise Level 4 (Solution)
\begin{frame}{In-Class Exercise --- Challenging Level (2/2)}
  \scriptsize
  We expand $X^2$ around the center of the narrow range: \pause

  \vspace{0.05cm}
  \begin{block}{Calculation}
    \scriptsize
    $$X^2=(X_0+\delta)^2=X_0^2+2X_0\delta+\delta^2\approx X_0^2+2X_0\delta \quad (\delta^2 \text{ negligible if } |\delta|\ll X_0)$$
  \end{block}

  \vspace{0.05cm}
  \pause
  \begin{alertblock}{Interpretation}
    $X^2$ becomes approximately linear in $X$ (slope $2X_0$), producing $r_{X,X^2}\approx1$ and a very high VIF. Predictions $\hat Y$ (which depend jointly on $c_1X+c_2X^2$) remain stable; what becomes difficult is separating the "effect of $X$" from the "effect of $X^2$" individually.
  \end{alertblock}
\end{frame}

% Slide 14: Section Closing
\begin{frame}{Section Closing: Nonlinear Transformations}
  \begin{reflexion}
    Polynomial regression demonstrates the flexibility of the least-squares framework: it suffices to expand the predictor space (adding $X^2$, $X^3$, or transformations like $\log X$) to model curvature, without ever abandoning the linear machinery derived in Section 09.07. The trade-off is the inherent multicollinearity between powers of the same variable.
  \end{reflexion}

  \vspace{0.2cm}
  \pause
  \begin{block}{Key Takeaways}
    \begin{itemize}\itemsep=0.08cm
      \item \textbf{Linearity in the parameters:} the true condition for OLS, not the shape of the relationship in $X$.
      \item \textbf{Partial F-Test:} formalizes whether a higher-order term is truly necessary.
      \item \textbf{Polynomial multicollinearity:} a high VIF between $X$ and $X^2$ does not invalidate predictions, only isolated interpretation.
    \end{itemize}
  \end{block}
\end{frame}

% Slide 15: Chapter Closing
\begin{frame}{Chapter 09 Closing: Linear and Multiple Regressions}
  \scriptsize
  \begin{retoclase}
    With this section we close the full Chapter 09: we began with \textbf{correlation} (09.01) as a conceptual premise, built \textbf{simple regression} from its motivation through its professional implementation (09.02-09.06), generalized to \textbf{multiple regression} via matrix notation and regularization (09.07), learned to \textbf{validate} (09.08) and \textbf{audit} (09.09) a model, implemented it with \textbf{industry tools} (09.10), and extended its reach to \textbf{categorical} predictors (09.11) and \textbf{nonlinear} relationships (09.12).
  \end{retoclase}

  \vspace{0.08cm}
  \pause
  \begin{alertblock}{The Common Thread}
    A single principle --- minimizing the sum of squared residuals --- unifies each of these twelve sections. Chapter 09 shows that linear regression, far from being a limited technique, is an extraordinarily flexible framework on which nearly all modern predictive methods in data science are built.
  \end{alertblock}
\end{frame}

\end{document}
